\chapter{Expected Outcomes}
\label{ch:outcomes}

Each research question can come out in more than one way.
Table~\ref{tab:out-map} fixes in advance how each outcome will be read, so that
negative results can be read as clearly as positive ones.

Two results hold whatever the data show. The cost breakdown of RQ1 is a result
in any case. And under a cold handover with an unchanged strong budget, the
repair rate cannot fall (Proposition~\ref{prop:nolost}). The open question of
RQ3 is therefore the size and sign of the saving, not a loss of repairs.

The expected picture is a saving of tens of percent in money and tokens, a
small change in time, and an open answer on the handover. Escalation is
expected to help most where the price gap is wide and the cheap tier repairs a
fair share of the bugs (Section~\ref{sec:sys-regimes}). The thesis reports
which of these conditions hold on the benchmark rather than assuming them.

{\footnotesize
\begin{longtable}{@{}>{\raggedright\arraybackslash}p{0.07\textwidth} >{\raggedright\arraybackslash}p{0.34\textwidth} >{\raggedright\arraybackslash}p{0.53\textwidth}@{}}
\caption{Possible outcomes and how they will be read.}
\label{tab:out-map}\\
\toprule
RQ & If \dots & then \dots\\
\midrule
\endfirsthead
\toprule
RQ & If \dots & then \dots\\
\midrule
\endhead
\bottomrule
\endlastfoot
RQ1 & the strong tier's calls dominate the cost & escalation is the main lever; its saving is bounded by the strong tier's spend on bugs the cheap tier repairs\\
RQ1 & validation dominates, even on the warm substrate & the tests are the cost; selection and the schedule matter more than the choice of model, and the saving of escalation is reported as diluted\\
RQ1 & most of the spend falls on bugs that are never repaired & no choice of model removes it; the Budget Guard is the lever, and its trade-off is reported\\
\midrule
RQ2 & M0 removes most of the build and start-up & cheap candidates cost little to check, and cheap-first costs little time\\
RQ2 & M0 disagrees with cold validation on some candidates & the rate and its causes are reported; the cold path is used wherever the verdict cannot be guaranteed\\
RQ2 & many bugs fail the calibration, or many candidates force a new runtime & the fallback rate bounds the saving, and those validations are charged at the cold price or with the new start\\
RQ2 & the schedule leaves much to a better order & the room is reported as the bound on what a new ordering could add; it is not pursued here\\
\midrule
RQ3 & \Step\ costs less than Strong-only & the bot meets the need; the breakdown shows where the saving came from\\
RQ3 & the cheap tier repairs too few bugs to pay & the break-even margin is negative and says by how much; for this setting the recommendation is to start at the strong tier\\
RQ3 & \Step\ is no better than Fixed-5 or Fixed-10 & the simpler rule is recommended; the stall signal adds nothing here\\
RQ3 & \Step\ is close to Hindsight & little is left for any router in test-driven repair, learned ones included\\
RQ3 & \Step\ is far from Hindsight & the gap is the room for a better trigger; whether it comes from hopeless bugs kept too long or from easy bugs given up too early is reported\\
RQ3 & money falls but time rises & the two prices are reported separately; which one matters is a property of the deployment\\
RQ3 & with equal attempts, the repair rate falls & the cheap attempts were worth less than the strong rounds they replaced; the budget split at which the loss disappears is reported\\
\midrule
RQ4 & a note lowers the strong tier's rounds & verified evidence transfers; the note is worth its tokens and becomes the default\\
RQ4 & no note differs from Cold & evidence does not transfer for this pair; cold start is recommended, in agreement with the fresh start of SWE-Router~\citep{son2026swerouter}\\
RQ4 & a note raises the anchoring rate or lowers repairs & the bias reported for drafts~\citep{ning2026revision} extends to this handover; the variant without patches says whether the drafts or the facts caused it\\
RQ4 & Transcript does as well as the notes & the summary loses nothing, and the shorter note is preferred for its tokens\\
RQ4 & the strong tier repairs nearly every bug in its first round & there is no room to measure a handover; RQ4 is reported on the pair one step down\\
\midrule
RQ5 & the bot costs less at the same repair rate & the system meets the need on this benchmark\\
RQ5 & the saving comes from the validation layer, not from escalation & the two prices of Strong-only separate them; the layer is recommended on its own, and escalation only where its break-even holds\\
RQ5 & no saving end to end & the saving is diluted by costs the system does not touch, mainly the strong tier's spend on bugs nobody repairs; the breakdown says which\\
RQ5 & cheap-tier patches match the developer's fix less often & plausible is not correct, and the cheap tier may overfit more; the saving is reported next to this gap\\
\end{longtable}
}
